% EMBASAMENTO TEÓRICO----------------------------------------------------------

\chapter{EMBASAMENTO TÉORICO}
\label{chap:embasamento-teorico}
Este capítulo apresenta, de forma clara, os conceitos e as tecnologias que fundamentam as atividades descritas neste relatório. Como o estágio teve por foco o desenvolvimento de um módulo de solicitações com a rastreabilidade, auditoria e reversão para o sistema Banco de Especialistas, é dada ênfase ao \textit{back-end} (Flask, SQLAlchemy e PostgreSQL) e aos conceitos de auditoria e versionamento de dados. As demais tecnologias (Vue.js, Nuxt, Pytest, Git e Docker) são abordadas como elementos de apoio ao desenvolvimento, à validação e à execução do módulo.

\section{Arquitetura Web Cliente-Servidor}
A arquitetura cliente-servidor organiza um sistema distribuído em dois papéis complementares: o cliente, que solicita serviços, e o servidor, que os provê. Na Web, o cliente é tipicamente um navegador que envia requisições por meio do protocolo HTTP, e o servidor processa essas requisições, aplica as regras de negócio, acessa a camada de persistência e devolve uma resposta \cite{Tanenbaum2007}.